% !TeX root = ../Book1.tex
\begin{figure}[htbp]
\centering
\begin{tikzpicture}[scale=1.1]
  \coordinate (v) at (0,0);
  \foreach \i in {1,...,5} {
    \coordinate (w\i) at ({90+72*(\i-1)}:2.2);
    \coordinate (c\i) at ({126+72*(\i-1)}:1.186);
  }
  \fill[color=AlgebraGray,opacity=0.08] (v)--(w5)--(w1)--cycle;
  \fill[color=AlgebraBlue,opacity=0.09] (v)--(w1)--(w2)--cycle;
  \draw[gray] (w1)--(w2)--(w3)--(w4)--(w5)--cycle;
  \foreach \i in {1,...,5} {\draw[gray] (v)--(w\i);}
  \draw[very thick,color=AlgebraBlue] (c1)--(c2)--(c3)--(c4)--(c5)--cycle;
  \draw[line width=1.9pt,color=AlgebraBlue] (c5)--(c1);
  \draw[line width=1.3pt] (v)--(w1);
  \foreach \i in {1,...,5} {
    \fill (w\i) circle (1.35pt);
    \fill[color=AlgebraBlue] (c\i) circle (1.5pt);
  }
  \fill (v) circle (1.7pt);
  \node[below] at (v) {\(v\)};
  \node[above] at (w1) {\(w_1\)};
  \node[left] at (w2) {\(w_2\)};
  \node[below left] at (w3) {\(w_3\)};
  \node[below right] at (w4) {\(w_4\)};
  \node[right] at (w5) {\(w_5\)};
  \node at (126:1.78) {\(F_1\)};
  \node at (54:1.78) {\(F_5\)};
  \node[left,color=AlgebraBlue] at (c1) {\(c_{F_1}\)};
  \node[right,color=AlgebraBlue] at (c5) {\(c_{F_5}\)};
  \node[color=AlgebraBlue] at (0,-1.61) {\(Q_v\)};
\end{tikzpicture}
\caption[正二十面体的邻面与面中心]{沿原顶点 \(v\) 的方向观察的局部投影。五个邻面 \(F_i=\operatorname{conv}\{v,w_i,w_{i+1}\}\)（下标按模 \(5\) 计算）的中心组成蓝色五边形 \(Q_v\)。原边 \([v,w_1]\) 邻接 \(F_5,F_1\)，对应加粗的边 \([c_{F_5},c_{F_1}]\)。}
\label{b1:fig:icosahedron-face-centers}
\end{figure}
